\section{A linear binomial radius and a signed radix in one residue class}
\label{sec:r20-direct}

All equations in this section are ordinary quartics over independently
unrestricted integers.  They use the numerical C222 proof calculus directly,
not a universal-polynomial machine-index substitution.  Relative to the
incoming factorial-free construction, two semantic edits eliminate two
coordinates and two residuals:
\[
 R=F^2\quad\longrightarrow\quad R=F,
 \qquad
 q=(2z+1)^2\quad\longrightarrow\quad q=4z+1,
 \qquad F=[2(t+1)]^t,\quad t=16+h^2.
\]
Here the variable power describes the proved meaning of an unchanged
quadratic Pell certificate; it is not an operation in the final quartic.
The first edit uses a classical lcm estimate.  The second allows an
initially negative radix and then proves its positivity in the necessary
order.  Exact residual-basis changes subsequently improve the collected
coefficients without changing the zero set.

Each final equation has the form
\[
 P(\boldsymbol x)=\sum_{i=0}^{m-1}R_i(\boldsymbol x)^2=0,
 \qquad\boldsymbol x\in\mathbb Z^n,
\]
with all $R_i$ of degree at most two and $P$ of exact degree four.  Every
residual and every collected coefficient is supplied, together with
coordinate and inverse row certificates.  The accompanying endpoint table
records the measured tradeoffs.  ``Height'' always means the largest
absolute coefficient after complete collection of $P$.
\begin{center}
\begin{tabular}{lrrrr}
\toprule
Profile & Variables & Rows & Monomials & Height\\
\midrule
Dimension & 239 & 175 & 12,658 & 6,144\\
Balanced & 240 & 176 & 5,232 & 4,096\\
Sparse & 240 & 176 & 4,927 & 7,596\\
\bottomrule
\end{tabular}
\end{center}
The 239-coordinate profile improves variables, residuals and collected
support over the incoming Shared Factorial Error 240-coordinate endpoint,
while retaining its height 6,144.  All three new profiles have maximum
residual coefficient 64.  These are achieved tradeoffs, not global minima.

\subsection{A classical lcm estimate makes the radius linear}
Write $\Lambda_n=\operatorname{lcm}(1,\ldots,n)$.  We recall the classical
bound $\Lambda_n\le4^n$, including a short proof so that the improvement
does not depend on an unexplained asymptotic estimate.  For $n\ge2$, set
$k=\lceil n/2\rceil$.  Then
\begin{equation}\label{eq:r20-lcm-recur}
 \Lambda_n\mid\Lambda_k\binom nk.
\end{equation}
Indeed, let $p^a$ be the largest power of $p$ not exceeding $n$.
If $p^a\le k$, it already divides $\Lambda_k$.  Otherwise
$p^{a-1}\le k$, so only one further factor $p$ is needed.
The $p^a$ term of Legendre's valuation formula for $\binom nk$ is one:
$n$ has one multiple of $p^a$, while $k$ and $n-k$ have none.
All other terms of that formula are nonnegative.

For $n=2m$, the middle binomial coefficient is at most $4^m$.
For $n=2m+1$, its two equal middle coefficients have sum at most
$2^{2m+1}$, so each is again at most $4^m$.  Complete induction in
\eqref{eq:r20-lcm-recur}, starting with $\Lambda_1=1$, yields
$\Lambda_n\le4^n$ in both cases.  This is a classical theorem, not a
claimed new number-theoretic result; see Chan and Norrish,
\emph{Proof Pearl: Bounding Least Common Multiples with Triangles},
J.~Automated Reasoning 62 (2019), 171--192, Theorem~3,
\href{https://doi.org/10.1007/s10817-017-9438-0}{doi:10.1007/s10817-017-9438-0}.

The unchanged horizon satisfies $t\ge16$.  Consequently
\begin{equation}\label{eq:r20-linear-radius}
 F=[2(t+1)]^t\ge34^t>16^t\ge\Lambda_t^2.
\end{equation}
The authoritative factorial-free predecessor, its
\texttt{report/main.tex}, proves the following two-sided binomial selector
lemma under the hypothesis $R\ge\Lambda_t^2$:
for $B=\binom{R+t}{t}$, every row $i\in\{1,\ldots,t\}$ has a prime power
$p^e$ exactly dividing $B$ with $p^e>t$ and $p^e\mid R+i$; conversely,
every prime power exactly dividing $B$ divides $R+i$ for some row.
The first assertion recovers every local proof row.  The second handles
\emph{all} primes in the reverse CRT construction, including primes not
selected to isolate a row.  Equation~\eqref{eq:r20-linear-radius} therefore
permits $R=F$.  An older predecessor review proves only a factorial-square
threshold; that weaker review is not the theorem used here.

\subsection{What is retained before signs are recovered}
The complete signed squared frame, semantic anchor, pooled Pell orientation
conditions, and all eighteen power slots are retained.  Both coordinates
of every principal Pell pair are oriented at their signed requested
indices.  No representative subset or frame atom is discarded.
The inherited signed-power lemma, using the strict frame bound
$|A|>W^W$, states that at a nonzero integer base $b$ with $|b|\ge2$, a
requested index $k$ must be nonnegative and its output equals $b^k$.
For $b=\pm1$, the output is $1$ or $(-1)^k$ even for signed $k$.
For a negative index at a nonunit base, the same small-error congruence
would force an integer identity $m b^{-k}=1$, which is impossible.
We never apply this lemma to the base zero.

The known positive index $t$ and base $2(t+1)$ first give
$F=[2(t+1)]^t>2^{32}$, independently of any radix sign.  The retained
beta certificate
\[
 v_\beta=(2w+1)^2,\qquad y=Fv_\beta
\]
then yields $y\ge F>0$.  The possible boundary $y=F$ is not silently
excluded; it is addressed after genuine local-row recovery below.

\subsection{The mod-four radix and the acyclic bootstrap}
The replacement $q=4z+1$ excludes both $q=0$ and $q=-1$.
It is stronger than oddness alone in precisely the needed respect:
$q=-1$ would make the numerator base $q+1$ zero.
We now exclude the only remaining unit, $q=1$.
In that case the base $q+1=2$ makes all nine actual CRT representative
exponents nonnegative.  Hence
\[
 H=2+\sum_\nu a_\nu\ge2,
 \qquad C_0\ge0.
\]
The retained raw index equation $B\mu=C_0+R+1$ has positive right-hand
side, so $B\ne0$ and $|BH|\ge2$.
The self-radix request gives $x\ge0$ and $1=(BH)^x$, forcing $x=0$.
But the last beta endpoint is
\[
 x=1+[1+(t+1)y]u_t.
\]
At $x=0$ it would make the integer $1+(t+1)y>1$ divide $-1$.
This contradiction uses neither a binomial interpretation nor a remainder
bound.

For every remaining radix, put $Q=|q|$.  Then $Q\ge3$ and
$|q+1|\ge2$.  The same argument recovers nonnegative representatives,
$H\ge2$, and $B\ne0$.  The derived label is another requested exponent
at $q+1$, so it is nonnegative; the unchanged anchor excludes zero.
The genuine self-radix equality is $q=(BH)^x$ with $x\ge0$.
Since $x=0$ would give $q=1$, we have $x\ge1$.
If the first beta quotient were $u_0\le-1$, its endpoint
\[
 x=F+1+(y+1)u_0
\]
would imply $x\le F-y\le0$.  Thus $u_0\ge0$, $x\ge F+1$, and the
last endpoint has $u_t\ge1$.  It follows that
\begin{equation}\label{eq:r20-q-size}
 x\ge2+(t+1)y\ge2+(t+1)F,\qquad
 Q=|BH|^x\ge2^x,\quad Q\ge4B^2,\quad |B|\le Q/4.
\end{equation}
No sign of $q$ or $B$ has yet been assumed.

All powers at $q$ and $q+1$ now have their genuine meanings at their
nonnegative indices.  In particular $M=q^x\ne0$ and every digit
denominator $L_i$ is a nonzero power of $q$; its exponent is one of
$t,t+1,y+1,x$, all at most $x$.
For $L_i v_i=M$ and $v_i U_i=M Q_i+R_i$, the nonzero integer $v_i$
divides $R_i$.  Hence $r_i=R_i/v_i$ is an integer and
$U_i=L_iQ_i+r_i$.  The unchanged common norm
\[
 M^2-1-4\sum_i R_i^2=a^2+b^2+c(c+1)
\]
gives $|r_i|<|L_i|/2$ without assuming signs for $L_i,v_i,M$.
Any congruent canonical tail of magnitude $<|L_i|/32$ must therefore
coincide with $r_i$: their difference is a multiple of $L_i$ of
magnitude less than $(1/2+1/32)|L_i|<|L_i|$.

\subsection{Authenticating the binomial at a signed radix}
Put $n=R+t$ and $C_j=\binom nj$.  The numerator request is exactly
$U=(q+1)^n$.  For $0\le j\le t$,
\[
 C_j\le(R+1)^t\le(F^2+1)^t=:V.
\]
The standard uniform binomial bound follows by Pascal induction and the
last-edge product bound.  Here $R=F\ge1$; the mod-four-only controls
instead retain $R=F^2$ and satisfy the same upper bound.
The symbol $V$ is proof notation, not a new power slot.
For $F\ge7$, $F^2+1\le2^{F-1}$ (the base case and induction are immediate).
Equation~\eqref{eq:r20-q-size} gives
\[
 V/Q\le2^{-F-t-2},\qquad V/(Q-1)<1/32.
\]
For either sign of $q$, the signed low polynomial
\[
 r^0=\sum_{j=0}^{t-1}C_jq^j
\]
has $|r^0|<Q^t/32$ and is congruent to $U$ modulo $q^t$.
The short-lift argument identifies the digit remainder with $r^0$.
Dividing the digit equation by $q^t$ and reducing modulo $q$ now gives
$B\equiv C_t\pmod q$.
But $0<C_t<Q/32$ and $|B|\le Q/4$, so $|B-C_t|<Q$:
\[
 B=\binom{R+t}{t}>0.
\]
Since $H>0$, the self-radix equality then proves $q=(BH)^x>0$.
The initially negative branch has been excluded by arithmetic.
The request $N=B^x$ now has its ordinary positive nonunit base meaning.
Also $B\ge R+t>R+1$, so the raw index equation excludes $C_0=0$.

The remaining range and local-congruence arguments are unchanged.
Their canonical tails are $<L_i/32$; all fixed front indices and local
majorants are no larger after replacing $F^2$ by $F$.
The strict endpoint budget $x\ge2+(t+1)F$ is preserved, including the
derived-label bound.  Applying the two-sided binomial selector theorem
using~\eqref{eq:r20-linear-radius} recovers every local C222 row.

\subsection{Recovering the proof and constructing all witnesses}
The recovered beta rows give accumulators $S_0,\ldots,S_t$ and updates
\[
 S_{i+1}(1+c_iF)=S_i(1+p_iF)(1+q_iF),
\]
whose weights are positive.  If $y=F$, the first endpoint would give
$S_0=0$, while the last gives $S_t=1$; every update propagates zero,
a contradiction.  Thus $y>F$, the initial remainder is genuinely
$F+1$, the final one is $1$, and all accumulators are positive.
Multiplication of the updates yields the inherited label-product identity.
The two sides have equal nonunit factor counts $d\le t$: if they differed,
the side with more factors would be at least $(F+1)^{d+1}$, while the
other is at most $t^d(F+1)^d$, contradicting $F+1>t^t$.
Each side is consequently the value at $F$ of a product of at most $t$
factors $1+jX$, $1\le j\le t$.  Its coefficient sum is at most
$(t+1)^t<F$.  Reduction modulo $F$ at the least nonzero degree shows
that a nonzero difference cannot vanish at $F$.
The products are identical polynomials; their distinct roots $-1/j$
recover every label multiplicity.  The retained Horn-closure argument
therefore gives a finite C222 refutation and hence a finite ZFC refutation.
This proves soundness.

Conversely, begin with a finite C222 refutation and choose a sufficiently
large padded horizon $t=16+h^2$.  Set $F=[2(t+1)]^t$, $R=F$, and
$B=\binom{F+t}{t}$.  The two-sided selector theorem applies by
\eqref{eq:r20-linear-radius}.
Let $O$ be the odd part of $t!$.  Choose an arbitrarily large odd multiple
$a_\beta$ of $O$, and put $y=F a_\beta^2$.
Since $v_2(F)\ge t>v_2(t!)$, this is an even multiple of $t!$ and can
exceed $F$ and every accumulator.  The beta CRT works as before; its
moduli are odd, so choose a sufficiently large even beta code $x$.

The smaller-radius binomial is still odd:
$v_2(F+i)=v_2(i)$ for $1\le i\le t$, since
$v_2(i)<t\le v_2(F)$.  Thus $v_2(B)=0$ and $N^2=B^{2x}$ is odd.
The all-primes CRT construction provides the actual representatives.
Adding $N^2$ permits their parities to be chosen while keeping the old
bounds.  Choose the dominant $a$ odd and every other representative even;
then $a,b$ have the required opposite parities, the nonzero semantic
anchor is even, and $H=2+\sum a_\nu$ is odd.
The changed front index satisfies $R+t\le B<B^2<N^2$, so every retained
principal-index and frame bound remains valid.
Now $q=(BH)^x$ is an odd square, hence $q\equiv1\pmod4$, and
$z=(q-1)/4$ supplies the new affine radix constraint.
Canonical normalized tails are $R_i=(M/L_i)r_i^0$ with
$0\le R_i<M/32$; the strict common norm gap has the usual ternary
representation.  The unchanged Pell completion supplies all other
witnesses.  Factorials are used only to select a divisible beta parameter,
not as an authenticated unknown.  This proves completeness.

\subsection{Literal source edits, exact inverse certificates, and limits}
In both pinned source equations $F=x_3$, $R=x_4$, and row zero is
$R-F^2$.  The source-interface audit checks that all previous factoring,
projection, translation, and row-basis steps retain this row and these
coordinates without using the row as a premise in an elimination.
Replacing it by $R-F$ therefore implements exactly the semantic radius
change.  The centered stored radix has actual value
$q=\text{stored}_q-1$; its old row is
$\text{stored}_q-2-4z-4z^2$.  The new row is
$\text{stored}_q-2-4z$, so its affine restoration is
$\text{stored}_q=4z+2$, not $4z+1$.
The second restoration is $R=F$.

These semantic edits preserve \emph{zero existence}, by the two ordinary
arithmetic directions above.  They are not claimed to be bijections
between all predecessor witness tuples.  Only after the source edits do
the two monic affine eliminations give exact restriction/restoration
bijections.  Every later row operation $R_i\leftarrow R_i+kR_j$ has the
integer inverse with $-k$ and preserves the common zero set, though it
generally changes the sum-of-squares polynomial.

The independent checkers import no new producer, verify all restored
coordinates and pulled-back residuals, reverse every frozen row operation,
collect every quartic coefficient, and check degree, coordinate liveness,
and all reported metrics.  Deliberate corruptions must fail, under both
normal Python and Python~\texttt{-O}.  These are finite artifact checks,
not replacements for the unbounded arithmetic proof.
The complete immediate predecessor, its ordinary proofs and calibration
interface sources, is retained under \texttt{inherited\_factorial\_free}.
The named logical interfaces are
\begin{quote}\small\ttfamily
KEndpoint.c222\_iff\_zf\_refutation\\
ConstructibleChoice.zf\_refutation\_iff\_zfc\_refutation.
\end{quote}
Their composition yields the stated equivalence with $\operatorname{Con}(\mathrm{ZFC})$.
No global minimum, proof of consistency, new proof-assistant closure, or
literature-wide priority is claimed.
